\subsection{All reflection images can fail the canonical-prime condition}

\paragraph{Finding A41 (a fixed-field simultaneous failure of the sufficient detector).}
Let $E/\mathbf F_{128}$ be the ordinary curve of A13 and A38.
There is a smooth genus-two curve $C/\mathbf F_{128}$ and a
constant etale cyclic cover $\pi:D\to C$ of degree five,
with ordinary reflection quotient $E$, together with a split
rational point $P\in C(\mathbf F_{128})$, having the following
property. For every rational lift of $P$ to $D$ and every
one of the five dihedral reflection quotients, its image
$R$ has
\[
                  \operatorname{ord}([K_E-2R])=4159,
             \qquad \operatorname{ord}_{4159}(2)=2079
                         \quad\hbox{odd}.
 \tag{AI.1}
\]
Here the quotient is identified with $E$ by the maps of A40;
the order assertion is independent of that identification.
Thus no choice of lift or reflection within this same
degree-five cover supplies the sufficient canonical-prime
hypothesis of A29. The assertion persists for these fixed
images after every extension of constants.

The proof supplies at least $95$ rational parameters for
the A40 inverse family over this field. Each resulting
cover has at least eleven distinct rational split source
points with the property just stated. These are existence
and counting statements in an explicit algebraic family;
no selected parameter or covering equation is claimed to
have been individually enumerated.

\begin{proof}
\emph{1. At most thirty-four parameters are excluded in A40.}
Use the notation of A40 over an algebraically closed field.
Let $F_1=\{P_0\}\times E$ and $F_2=E\times\{Q_0\}$ be
the two fibre classes on $E\times E$. The critical curve
\[
 Z=\{(P,Q):\gamma(Q)=\overline T\gamma(P)\},
         \qquad h(P,Q)=Tz_P-2z_Q
 \tag{AI.2}
\]
is the pullback of the diagonal in $\mathbf P^1\times\mathbf P^1$.
Both canonical maps have degree two, so its numerical
divisor class is
\[
                             [Z]=2F_1+2F_2.
 \tag{AI.3}
\]
This is also the class of its reduced curve used in A40.
Indeed its projections are generically separable, so the
Cartier divisor is generically reduced; a Cartier divisor
on this smooth surface is Cohen--Macaulay and has no
embedded associated components. It is therefore reduced.

Put $H=h^*\mathcal O_J(T\Theta)$. Its restrictions to the
two fibres have degrees
\[
                         \deg(H|_{F_1})=12,
                \qquad \deg(H|_{F_2})=2.
 \tag{AI.4}
\]
For the first equality, varying $Q$ gives a translate of
$-[2]\Theta$, so pullback by $[2]$ multiplies the divisor
class by four \cite[Corollary~6.6 and Remark~12.13,
pp.~9 and~24]{MilneAV}; (ID.4) then gives
$4(\Theta\cdot T\Theta)=12$.
For the second, varying $P$ gives a translate of $T\Theta$,
whose self-intersection is two. Translation does not alter
these degrees. Hence
\[
                             Z\cdot H=28.
 \tag{AI.5}
\]

Step~3 of A40 proves that no component of $Z$ maps into
$T\Theta$ under $h$. The pulled-back effective divisor
therefore meets $Z$ properly, and the sum of its local
intersection lengths is exactly $28$. In particular there
are at most $28$ geometric intersection points. Each
off-vertical singularity of the residual curve $D_t$ gives
one of these points with $h(P,Q)=Tz_t$, and $Tz_t$
determines $t$ uniquely. Thus at most $28$ parameters can
produce such a singularity.

The other excluded parameters in A40 satisfy
$\overline T\gamma(t)=\gamma(W)$ for a Weierstrass point
$W$. There are three branch values, each with at most two
preimages under the degree-two map $\overline T\gamma$.
This excludes at most six additional parameters. Consequently
\[
        \#\{t\in E(\bar k):D_t\hbox{ is not smooth}\}\le34.
 \tag{AI.6}
\]
The integrality of every residual curve was already proved
in A40, so no additional reducibility exclusions are needed.

\emph{2. The inverse covers already exist over $\mathbf F_{128}$.}
For the A13 curve, choose $B=\mathcal O_E(W)$ at one of
its three $\mathbf F_2$-rational Weierstrass points. This
is a theta characteristic and makes the symmetric theta
embedding in A40 defined over $\mathbf F_2$.
The integral Rosati endomorphism $T$ is also defined over
$\mathbf F_2$, by the checked supplement to A33.
Thus the entire family $D_t$, its two projection maps,
and its dihedral data are defined over the field of $t$.
Finding A38 gives $\#E(\mathbf F_{128})=129$.
By (AI.6), at least
\[
                              129-34=95
 \tag{AI.7}
\]
of these rational parameters give smooth genus-six curves
$D_t$ over $\mathbf F_{128}$. For any such $t$, A40 gives
a genus-two quotient $C_t$ and a constant etale cyclic
degree-five map over this same field. No further extension
is required to define its deck group or reflections.

The pullback, norm, and complementary-factor construction
of A40 are all defined over $\mathbf F_{128}$. They give
the isogeny over that field
\[
                      J_{D_t}\sim J_E^2\times J_{C_t}.
 \tag{AI.8}
\]
The Frobenius trace of $E$ over this field is zero, again
by A38. Taking traces in (AI.8) and applying the curve
point-count formula therefore gives
\[
 \begin{gathered}
 \#D_t(\mathbf F_{128})=\#C_t(\mathbf F_{128})
       \ge129-4\sqrt{128}>83,\\
                    \#D_t(\mathbf F_{128})\ge84.
 \end{gathered}
 \tag{AI.9}
\]
Here the inequality is the genus-two Weil bound
\cite[Theorem~11.1]{MilneJV}. It applies to $C_t$ and
does not require an explicit equation or a point-count
computation for that individual curve.

\emph{3. All five images can be simultaneously non-Weierstrass.}
Write $v:D_t\to E$ for its first projection and let
$\sigma$ generate its constant degree-five deck group.
The five maps
\[
                         v\sigma^j:D_t\longrightarrow E,
                              \qquad 0\le j<5,
 \tag{AI.10}
\]
are the five reflection quotient maps, using the
$\mathbf F_{128}$-defined identifications with $E$.
Each has degree two. The union of the inverse images of
the three Weierstrass points under all five maps has at
most
\[
                               5\cdot3\cdot2=30
 \tag{AI.11}
\]
geometric points, and hence at most that many rational
points. By (AI.9), at least $54$ rational points $P'$ of
$D_t$ have all five images non-Weierstrass. For each of
them A38 proves (AI.1) for every image.

This good set is invariant under $\sigma$, since applying
$\sigma$ merely permutes the five maps. The degree-five
action is free. Thus its size, already at least $54$, is
a multiple of five and is at least $55$. It supplies at
least eleven distinct points $P=\pi(P')$ of $C_t$.
For each, the entire fibre $\{\sigma^jP'\}$ is rational,
so $P$ is split. Changing the lift permutes the same five
images. The five reflections are the conjugates
$\sigma^{-j}\tau\sigma^j$, and thus (AI.10) covers every
reflection choice. A different isomorphism of a quotient
with $E$ preserves the order of its canonical divisor
class. Finally, constant extension preserves the exact
order of each of these fixed classes. This proves every
assertion.
\end{proof}

\paragraph{The remaining choice of cover.}
A41 rules out forcing the A29 canonical-prime hypothesis
by changing only the lift or reflection of a fixed
quintic transfer. It leaves open the possibility of
choosing a different cyclic cover for the prescribed
source point, or using a geometric theta witness or a
different rational-Prym criterion. A43 below determines the
two-rank of these sources to be one and proves that both
Weierstrass points are rational and split with the same failure
at every reflection image. The order-$4159$
geometric theta witness on $E$ from A13 is compatible
with all these canonical failures, so no failure of
point-section separation follows.
